% Options for packages loaded elsewhere
\PassOptionsToPackage{unicode}{hyperref}
\PassOptionsToPackage{hyphens}{url}
\documentclass[
]{article}
\usepackage{xcolor}
\usepackage[margin=1in]{geometry}
\usepackage{amsmath,amssymb}
\setcounter{secnumdepth}{-\maxdimen} % remove section numbering
\usepackage{iftex}
\ifPDFTeX
  \usepackage[T1]{fontenc}
  \usepackage[utf8]{inputenc}
  \usepackage{textcomp} % provide euro and other symbols
\else % if luatex or xetex
  \usepackage{unicode-math} % this also loads fontspec
  \defaultfontfeatures{Scale=MatchLowercase}
  \defaultfontfeatures[\rmfamily]{Ligatures=TeX,Scale=1}
\fi
\usepackage{lmodern}
\ifPDFTeX\else
  % xetex/luatex font selection
\fi
% Use upquote if available, for straight quotes in verbatim environments
\IfFileExists{upquote.sty}{\usepackage{upquote}}{}
\IfFileExists{microtype.sty}{% use microtype if available
  \usepackage[]{microtype}
  \UseMicrotypeSet[protrusion]{basicmath} % disable protrusion for tt fonts
}{}
\makeatletter
\@ifundefined{KOMAClassName}{% if non-KOMA class
  \IfFileExists{parskip.sty}{%
    \usepackage{parskip}
  }{% else
    \setlength{\parindent}{0pt}
    \setlength{\parskip}{6pt plus 2pt minus 1pt}}
}{% if KOMA class
  \KOMAoptions{parskip=half}}
\makeatother
\usepackage{color}
\usepackage{fancyvrb}
\newcommand{\VerbBar}{|}
\newcommand{\VERB}{\Verb[commandchars=\\\{\}]}
\DefineVerbatimEnvironment{Highlighting}{Verbatim}{commandchars=\\\{\}}
% Add ',fontsize=\small' for more characters per line
\usepackage{framed}
\definecolor{shadecolor}{RGB}{248,248,248}
\newenvironment{Shaded}{\begin{snugshade}}{\end{snugshade}}
\newcommand{\AlertTok}[1]{\textcolor[rgb]{0.94,0.16,0.16}{#1}}
\newcommand{\AnnotationTok}[1]{\textcolor[rgb]{0.56,0.35,0.01}{\textbf{\textit{#1}}}}
\newcommand{\AttributeTok}[1]{\textcolor[rgb]{0.13,0.29,0.53}{#1}}
\newcommand{\BaseNTok}[1]{\textcolor[rgb]{0.00,0.00,0.81}{#1}}
\newcommand{\BuiltInTok}[1]{#1}
\newcommand{\CharTok}[1]{\textcolor[rgb]{0.31,0.60,0.02}{#1}}
\newcommand{\CommentTok}[1]{\textcolor[rgb]{0.56,0.35,0.01}{\textit{#1}}}
\newcommand{\CommentVarTok}[1]{\textcolor[rgb]{0.56,0.35,0.01}{\textbf{\textit{#1}}}}
\newcommand{\ConstantTok}[1]{\textcolor[rgb]{0.56,0.35,0.01}{#1}}
\newcommand{\ControlFlowTok}[1]{\textcolor[rgb]{0.13,0.29,0.53}{\textbf{#1}}}
\newcommand{\DataTypeTok}[1]{\textcolor[rgb]{0.13,0.29,0.53}{#1}}
\newcommand{\DecValTok}[1]{\textcolor[rgb]{0.00,0.00,0.81}{#1}}
\newcommand{\DocumentationTok}[1]{\textcolor[rgb]{0.56,0.35,0.01}{\textbf{\textit{#1}}}}
\newcommand{\ErrorTok}[1]{\textcolor[rgb]{0.64,0.00,0.00}{\textbf{#1}}}
\newcommand{\ExtensionTok}[1]{#1}
\newcommand{\FloatTok}[1]{\textcolor[rgb]{0.00,0.00,0.81}{#1}}
\newcommand{\FunctionTok}[1]{\textcolor[rgb]{0.13,0.29,0.53}{\textbf{#1}}}
\newcommand{\ImportTok}[1]{#1}
\newcommand{\InformationTok}[1]{\textcolor[rgb]{0.56,0.35,0.01}{\textbf{\textit{#1}}}}
\newcommand{\KeywordTok}[1]{\textcolor[rgb]{0.13,0.29,0.53}{\textbf{#1}}}
\newcommand{\NormalTok}[1]{#1}
\newcommand{\OperatorTok}[1]{\textcolor[rgb]{0.81,0.36,0.00}{\textbf{#1}}}
\newcommand{\OtherTok}[1]{\textcolor[rgb]{0.56,0.35,0.01}{#1}}
\newcommand{\PreprocessorTok}[1]{\textcolor[rgb]{0.56,0.35,0.01}{\textit{#1}}}
\newcommand{\RegionMarkerTok}[1]{#1}
\newcommand{\SpecialCharTok}[1]{\textcolor[rgb]{0.81,0.36,0.00}{\textbf{#1}}}
\newcommand{\SpecialStringTok}[1]{\textcolor[rgb]{0.31,0.60,0.02}{#1}}
\newcommand{\StringTok}[1]{\textcolor[rgb]{0.31,0.60,0.02}{#1}}
\newcommand{\VariableTok}[1]{\textcolor[rgb]{0.00,0.00,0.00}{#1}}
\newcommand{\VerbatimStringTok}[1]{\textcolor[rgb]{0.31,0.60,0.02}{#1}}
\newcommand{\WarningTok}[1]{\textcolor[rgb]{0.56,0.35,0.01}{\textbf{\textit{#1}}}}
\usepackage{graphicx}
\makeatletter
\newsavebox\pandoc@box
\newcommand*\pandocbounded[1]{% scales image to fit in text height/width
  \sbox\pandoc@box{#1}%
  \Gscale@div\@tempa{\textheight}{\dimexpr\ht\pandoc@box+\dp\pandoc@box\relax}%
  \Gscale@div\@tempb{\linewidth}{\wd\pandoc@box}%
  \ifdim\@tempb\p@<\@tempa\p@\let\@tempa\@tempb\fi% select the smaller of both
  \ifdim\@tempa\p@<\p@\scalebox{\@tempa}{\usebox\pandoc@box}%
  \else\usebox{\pandoc@box}%
  \fi%
}
% Set default figure placement to htbp
\def\fps@figure{htbp}
\makeatother
\setlength{\emergencystretch}{3em} % prevent overfull lines
\providecommand{\tightlist}{%
  \setlength{\itemsep}{0pt}\setlength{\parskip}{0pt}}
\usepackage{bookmark}
\IfFileExists{xurl.sty}{\usepackage{xurl}}{} % add URL line breaks if available
\urlstyle{same}
\hypersetup{
  pdftitle={Inference for numerical data},
  pdfauthor={Tenzin Thakuri},
  hidelinks,
  pdfcreator={LaTeX via pandoc}}

\title{Inference for numerical data}
\author{Tenzin Thakuri}
\date{}

\begin{document}
\maketitle

\subsection{Getting Started}\label{getting-started}

\subsubsection{Load packages}\label{load-packages}

In this lab, we will explore and visualize the data using the
\textbf{tidyverse} suite of packages, and perform statistical inference
using \textbf{infer}. The data can be found in the companion package for
OpenIntro resources, \textbf{openintro}.

Let's load the packages.

\begin{Shaded}
\begin{Highlighting}[]
\FunctionTok{library}\NormalTok{(tidyverse)}
\FunctionTok{library}\NormalTok{(openintro)}
\FunctionTok{library}\NormalTok{(infer)}
\end{Highlighting}
\end{Shaded}

\subsubsection{The data}\label{the-data}

Every two years, the Centers for Disease Control and Prevention conduct
the Youth Risk Behavior Surveillance System (YRBSS) survey, where it
takes data from high schoolers (9th through 12th grade), to analyze
health patterns. You will work with a selected group of variables from a
random sample of observations during one of the years the YRBSS was
conducted.

Load the \texttt{yrbss} data set into your workspace.

\begin{Shaded}
\begin{Highlighting}[]
\FunctionTok{data}\NormalTok{(}\StringTok{\textquotesingle{}yrbss\textquotesingle{}}\NormalTok{, }\AttributeTok{package=}\StringTok{\textquotesingle{}openintro\textquotesingle{}}\NormalTok{)}
\end{Highlighting}
\end{Shaded}

There are observations on 13 different variables, some categorical and
some numerical. The meaning of each variable can be found by bringing up
the help file:

\begin{Shaded}
\begin{Highlighting}[]
\NormalTok{?yrbss}
\end{Highlighting}
\end{Shaded}

\begin{enumerate}
\def\labelenumi{\arabic{enumi}.}
\tightlist
\item
  What are the cases in this data set? How many cases are there in our
  sample?
\end{enumerate}

\textbf{Insert your answer here}

The cases are the individual students who participated in the
survey.There are 13,583 cases in our sample.

Remember that you can answer this question by viewing the data in the
data viewer or by using the following command:

\begin{Shaded}
\begin{Highlighting}[]
\FunctionTok{glimpse}\NormalTok{(yrbss)}
\end{Highlighting}
\end{Shaded}

\begin{verbatim}
## Rows: 13,583
## Columns: 13
## $ age                      <int> 14, 14, 15, 15, 15, 15, 15, 14, 15, 15, 15, 1~
## $ gender                   <chr> "female", "female", "female", "female", "fema~
## $ grade                    <chr> "9", "9", "9", "9", "9", "9", "9", "9", "9", ~
## $ hispanic                 <chr> "not", "not", "hispanic", "not", "not", "not"~
## $ race                     <chr> "Black or African American", "Black or Africa~
## $ height                   <dbl> NA, NA, 1.73, 1.60, 1.50, 1.57, 1.65, 1.88, 1~
## $ weight                   <dbl> NA, NA, 84.37, 55.79, 46.72, 67.13, 131.54, 7~
## $ helmet_12m               <chr> "never", "never", "never", "never", "did not ~
## $ text_while_driving_30d   <chr> "0", NA, "30", "0", "did not drive", "did not~
## $ physically_active_7d     <int> 4, 2, 7, 0, 2, 1, 4, 4, 5, 0, 0, 0, 4, 7, 7, ~
## $ hours_tv_per_school_day  <chr> "5+", "5+", "5+", "2", "3", "5+", "5+", "5+",~
## $ strength_training_7d     <int> 0, 0, 0, 0, 1, 0, 2, 0, 3, 0, 3, 0, 0, 7, 7, ~
## $ school_night_hours_sleep <chr> "8", "6", "<5", "6", "9", "8", "9", "6", "<5"~
\end{verbatim}

\subsection{Exploratory data analysis}\label{exploratory-data-analysis}

You will first start with analyzing the weight of the participants in
kilograms: \texttt{weight}.

Using visualization and summary statistics, describe the distribution of
weights. The \texttt{summary} function can be useful.

\begin{Shaded}
\begin{Highlighting}[]
\FunctionTok{summary}\NormalTok{(yrbss}\SpecialCharTok{$}\NormalTok{weight)}
\end{Highlighting}
\end{Shaded}

\begin{verbatim}
##    Min. 1st Qu.  Median    Mean 3rd Qu.    Max.    NA's 
##   29.94   56.25   64.41   67.91   76.20  180.99    1004
\end{verbatim}

\begin{enumerate}
\def\labelenumi{\arabic{enumi}.}
\setcounter{enumi}{1}
\tightlist
\item
  How many observations are we missing weights from? The distribution
  shape is right-skewed. The mean weight is 67.91 which is higher than
  Median 64.41 The weight range from 29.94 to 180.99. The middle 50\% of
  weights fall between 56.25 (Q1) and 76.20 (Q3). We are missing 1004
  observations for weights in the dataset.
\end{enumerate}

\textbf{Insert your answer here}

\begin{Shaded}
\begin{Highlighting}[]
\FunctionTok{ggplot}\NormalTok{(yrbss, }\FunctionTok{aes}\NormalTok{(}\AttributeTok{x =}\NormalTok{ weight)) }\SpecialCharTok{+}
  \FunctionTok{geom\_histogram}\NormalTok{(}
    \AttributeTok{bins =} \DecValTok{30}\NormalTok{,}
    \AttributeTok{fill =} \StringTok{"steelblue"}\NormalTok{,}
    \AttributeTok{color =} \StringTok{"white"}
\NormalTok{  ) }\SpecialCharTok{+}
  \FunctionTok{labs}\NormalTok{(}
    \AttributeTok{title =} \StringTok{"Distribution of Student Weight"}\NormalTok{,}
    \AttributeTok{x =} \StringTok{"Weight"}\NormalTok{,}
    \AttributeTok{y =} \StringTok{"Number of Students"}
\NormalTok{  ) }\SpecialCharTok{+}
  \FunctionTok{theme\_minimal}\NormalTok{()}
\end{Highlighting}
\end{Shaded}

\pandocbounded{\includegraphics[keepaspectratio]{Lab7_inf_for_numerical_data_files/figure-latex/unnamed-chunk-1-1.pdf}}

Next, consider the possible relationship between a high schooler's
weight and their physical activity. Plotting the data is a useful first
step because it helps us quickly visualize trends, identify strong
associations, and develop research questions.

First, let's create a new variable \texttt{physical\_3plus}, which will
be coded as either ``yes'' if they are physically active for at least 3
days a week, and ``no'' if not.

\begin{Shaded}
\begin{Highlighting}[]
\NormalTok{yrbss }\OtherTok{\textless{}{-}}\NormalTok{ yrbss }\SpecialCharTok{\%\textgreater{}\%} 
  \FunctionTok{mutate}\NormalTok{(}\AttributeTok{physical\_3plus =} \FunctionTok{ifelse}\NormalTok{(yrbss}\SpecialCharTok{$}\NormalTok{physically\_active\_7d }\SpecialCharTok{\textgreater{}} \DecValTok{2}\NormalTok{, }\StringTok{"yes"}\NormalTok{, }\StringTok{"no"}\NormalTok{))}
\end{Highlighting}
\end{Shaded}

\begin{enumerate}
\def\labelenumi{\arabic{enumi}.}
\setcounter{enumi}{2}
\tightlist
\item
  Make a side-by-side boxplot of \texttt{physical\_3plus} and
  \texttt{weight}. Is there a relationship between these two variables?
  What did you expect and why?
\end{enumerate}

\textbf{Insert your answer here}

\begin{Shaded}
\begin{Highlighting}[]
\FunctionTok{ggplot}\NormalTok{(yrbss, }\FunctionTok{aes}\NormalTok{(}\AttributeTok{x =}\NormalTok{ physical\_3plus, }\AttributeTok{y =}\NormalTok{ weight)) }\SpecialCharTok{+}
  \FunctionTok{geom\_boxplot}\NormalTok{(}\AttributeTok{fill =} \StringTok{"lightblue"}\NormalTok{) }\SpecialCharTok{+}
  \FunctionTok{labs}\NormalTok{(}
    \AttributeTok{title =} \StringTok{"Weight by Physical Activity Level"}\NormalTok{,}
    \AttributeTok{x =} \StringTok{"Physically Active 3 or More Days"}\NormalTok{,}
    \AttributeTok{y =} \StringTok{"Weight"}
\NormalTok{  ) }\SpecialCharTok{+}
  \FunctionTok{theme\_minimal}\NormalTok{()}
\end{Highlighting}
\end{Shaded}

\pandocbounded{\includegraphics[keepaspectratio]{Lab7_inf_for_numerical_data_files/figure-latex/unnamed-chunk-2-1.pdf}}
There appears to be a very weak relationship between physical\_3plus and
weight. Median weight of student physically active at least 3 days a
week is slightly higher than student who are not. I expected the weight
of student who is physically active at least 3 days a week is less than
the student who are not, because being physically active burns calories
and weight loss.

The box plots show how the medians of the two distributions compare, but
we can also compare the means of the distributions using the following
to first group the data by the \texttt{physical\_3plus} variable, and
then calculate the mean \texttt{weight} in these groups using the
\texttt{mean} function while ignoring missing values by setting the
\texttt{na.rm} argument to \texttt{TRUE}.

\begin{Shaded}
\begin{Highlighting}[]
\NormalTok{yrbss }\SpecialCharTok{\%\textgreater{}\%}
  \FunctionTok{group\_by}\NormalTok{(physical\_3plus) }\SpecialCharTok{\%\textgreater{}\%}
  \FunctionTok{summarise}\NormalTok{(}\AttributeTok{mean\_weight =} \FunctionTok{mean}\NormalTok{(weight, }\AttributeTok{na.rm =} \ConstantTok{TRUE}\NormalTok{))}
\end{Highlighting}
\end{Shaded}

\begin{verbatim}
## # A tibble: 3 x 2
##   physical_3plus mean_weight
##   <chr>                <dbl>
## 1 no                    66.7
## 2 yes                   68.4
## 3 <NA>                  69.9
\end{verbatim}

There is an observed difference, but is this difference statistically
significant? In order to answer this question we will conduct a
hypothesis test.

\subsection{Inference}\label{inference}

\begin{enumerate}
\def\labelenumi{\arabic{enumi}.}
\setcounter{enumi}{3}
\tightlist
\item
  Are all conditions necessary for inference satisfied? Comment on each.
  You can compute the group sizes with the \texttt{summarize} command
  above by defining a new variable with the definition \texttt{n()}.
\end{enumerate}

\textbf{Insert your answer here}

\begin{Shaded}
\begin{Highlighting}[]
\NormalTok{yrbss }\SpecialCharTok{\%\textgreater{}\%}
  \FunctionTok{group\_by}\NormalTok{(physical\_3plus) }\SpecialCharTok{\%\textgreater{}\%}
  \FunctionTok{summarise}\NormalTok{(}\AttributeTok{mean\_weight =} \FunctionTok{mean}\NormalTok{(weight, }\AttributeTok{na.rm =} \ConstantTok{TRUE}\NormalTok{),}\AttributeTok{n=}\FunctionTok{n}\NormalTok{())}
\end{Highlighting}
\end{Shaded}

\begin{verbatim}
## # A tibble: 3 x 3
##   physical_3plus mean_weight     n
##   <chr>                <dbl> <int>
## 1 no                    66.7  4404
## 2 yes                   68.4  8906
## 3 <NA>                  69.9   273
\end{verbatim}

Yes, all conditions are necessary for inference satisfied. Independence:
In this dataset one student weight does not influence another's.
Students who are physically active for at least 3 days is (yes 8906) and
who are not is (no 4404). So they are independent of each other.

Sample size and skew: The boxplot shows weight is right-skewed with many
high outliers in both groups. Since both condition are met, it is
appropriate to do inference (a confidence interval or hypothesis test).

\begin{enumerate}
\def\labelenumi{\arabic{enumi}.}
\setcounter{enumi}{4}
\tightlist
\item
  Write the hypotheses for testing if the average weights are different
  for those who exercise at least times a week and those who don't.
\end{enumerate}

\textbf{Insert your answer here} Let μ₁ be the true mean weight of
students who are physically active at least 3 days a week, and μ₂ the
true mean weight of students who are not.

Null hypothesis (H₀): The average weights are the same for the two
groups.

H₀: μ₁ = μ₂ (or μ₁ − μ₂ = 0)

Alternative hypothesis (Hₐ): The average weights are different for the
two groups.

Hₐ: μ₁ ≠ μ₂ (or μ₁ − μ₂ ≠ 0)

Next, we will introduce a new function, \texttt{hypothesize}, that falls
into the \texttt{infer} workflow. You will use this method for
conducting hypothesis tests.

But first, we need to initialize the test, which we will save as
\texttt{obs\_diff}.

\begin{Shaded}
\begin{Highlighting}[]
\NormalTok{obs\_diff }\OtherTok{\textless{}{-}}\NormalTok{ yrbss }\SpecialCharTok{\%\textgreater{}\%}
  \FunctionTok{drop\_na}\NormalTok{(physical\_3plus) }\SpecialCharTok{\%\textgreater{}\%}
  \FunctionTok{specify}\NormalTok{(weight }\SpecialCharTok{\textasciitilde{}}\NormalTok{ physical\_3plus) }\SpecialCharTok{\%\textgreater{}\%}
  \FunctionTok{calculate}\NormalTok{(}\AttributeTok{stat =} \StringTok{"diff in means"}\NormalTok{, }\AttributeTok{order =} \FunctionTok{c}\NormalTok{(}\StringTok{"yes"}\NormalTok{, }\StringTok{"no"}\NormalTok{))}
\end{Highlighting}
\end{Shaded}

Notice how you can use the functions \texttt{specify} and
\texttt{calculate} again like you did for calculating confidence
intervals. Here, though, the statistic you are searching for is the
difference in means, with the order being \texttt{yes\ -\ no\ !=\ 0}.

After you have initialized the test, you need to simulate the test on
the null distribution, which we will save as \texttt{null}.

\begin{Shaded}
\begin{Highlighting}[]
\NormalTok{null\_dist }\OtherTok{\textless{}{-}}\NormalTok{ yrbss }\SpecialCharTok{\%\textgreater{}\%}
  \FunctionTok{drop\_na}\NormalTok{(physical\_3plus) }\SpecialCharTok{\%\textgreater{}\%}
  \FunctionTok{specify}\NormalTok{(weight }\SpecialCharTok{\textasciitilde{}}\NormalTok{ physical\_3plus) }\SpecialCharTok{\%\textgreater{}\%}
  \FunctionTok{hypothesize}\NormalTok{(}\AttributeTok{null =} \StringTok{"independence"}\NormalTok{) }\SpecialCharTok{\%\textgreater{}\%}
  \FunctionTok{generate}\NormalTok{(}\AttributeTok{reps =} \DecValTok{1000}\NormalTok{, }\AttributeTok{type =} \StringTok{"permute"}\NormalTok{) }\SpecialCharTok{\%\textgreater{}\%}
  \FunctionTok{calculate}\NormalTok{(}\AttributeTok{stat =} \StringTok{"diff in means"}\NormalTok{, }\AttributeTok{order =} \FunctionTok{c}\NormalTok{(}\StringTok{"yes"}\NormalTok{, }\StringTok{"no"}\NormalTok{))}
\end{Highlighting}
\end{Shaded}

Here, \texttt{hypothesize} is used to set the null hypothesis as a test
for independence. In one sample cases, the \texttt{null} argument can be
set to ``point'' to test a hypothesis relative to a point estimate.

Also, note that the \texttt{type} argument within \texttt{generate} is
set to \texttt{permute}, whichis the argument when generating a null
distribution for a hypothesis test.

We can visualize this null distribution with the following code:

\begin{Shaded}
\begin{Highlighting}[]
\FunctionTok{ggplot}\NormalTok{(}\AttributeTok{data =}\NormalTok{ null\_dist, }\FunctionTok{aes}\NormalTok{(}\AttributeTok{x =}\NormalTok{ stat)) }\SpecialCharTok{+}
  \FunctionTok{geom\_histogram}\NormalTok{()}
\end{Highlighting}
\end{Shaded}

\pandocbounded{\includegraphics[keepaspectratio]{Lab7_inf_for_numerical_data_files/figure-latex/unnamed-chunk-4-1.pdf}}

\begin{enumerate}
\def\labelenumi{\arabic{enumi}.}
\setcounter{enumi}{5}
\tightlist
\item
  How many of these \texttt{null} permutations have a difference of at
  least \texttt{obs\_stat}?
\end{enumerate}

\textbf{Insert your answer here}

\begin{Shaded}
\begin{Highlighting}[]
\NormalTok{obs\_stat }\OtherTok{\textless{}{-}}\NormalTok{ obs\_diff }\SpecialCharTok{\%\textgreater{}\%}
  \FunctionTok{pull}\NormalTok{(stat)}

\NormalTok{obs\_stat}
\end{Highlighting}
\end{Shaded}

\begin{verbatim}
## [1] 1.774584
\end{verbatim}

Now that the test is initialized and the null distribution formed, you
can calculate the p-value for your hypothesis test using the function
\texttt{get\_p\_value}.

\begin{Shaded}
\begin{Highlighting}[]
\NormalTok{null\_dist }\SpecialCharTok{\%\textgreater{}\%}
  \FunctionTok{get\_p\_value}\NormalTok{(}\AttributeTok{obs\_stat =}\NormalTok{ obs\_diff, }\AttributeTok{direction =} \StringTok{"two\_sided"}\NormalTok{)}
\end{Highlighting}
\end{Shaded}

\begin{verbatim}
## # A tibble: 1 x 1
##   p_value
##     <dbl>
## 1       0
\end{verbatim}

\begin{Shaded}
\begin{Highlighting}[]
\FunctionTok{sum}\NormalTok{(}\FunctionTok{abs}\NormalTok{(null\_dist}\SpecialCharTok{$}\NormalTok{stat) }\SpecialCharTok{\textgreater{}=} \FunctionTok{abs}\NormalTok{(obs\_stat))}
\end{Highlighting}
\end{Shaded}

\begin{verbatim}
## [1] 0
\end{verbatim}

\begin{Shaded}
\begin{Highlighting}[]
\FunctionTok{sum}\NormalTok{(}\FunctionTok{abs}\NormalTok{(null\_dist}\SpecialCharTok{$}\NormalTok{stat) }\SpecialCharTok{\textgreater{}=} \FunctionTok{abs}\NormalTok{(obs\_stat)) }\SpecialCharTok{/} \DecValTok{1000}
\end{Highlighting}
\end{Shaded}

\begin{verbatim}
## [1] 0
\end{verbatim}

The number of null permutations with a difference at least as extreme as
the observed statistic is 0 which means none of the 1,000 simulated null
differences were as extreme as the observed difference This the standard
workflow for performing hypothesis tests.

\begin{enumerate}
\def\labelenumi{\arabic{enumi}.}
\setcounter{enumi}{6}
\tightlist
\item
  Construct and record a confidence interval for the difference between
  the weights of those who exercise at least three times a week and
  those who don't, and interpret this interval in context of the data.
\end{enumerate}

\begin{center}\rule{0.5\linewidth}{0.5pt}\end{center}

\begin{Shaded}
\begin{Highlighting}[]
\FunctionTok{set.seed}\NormalTok{(}\DecValTok{123}\NormalTok{)}

\NormalTok{weight\_ci }\OtherTok{\textless{}{-}}\NormalTok{ yrbss }\SpecialCharTok{\%\textgreater{}\%}
  \FunctionTok{drop\_na}\NormalTok{(weight, physical\_3plus) }\SpecialCharTok{\%\textgreater{}\%}
  \FunctionTok{specify}\NormalTok{(weight }\SpecialCharTok{\textasciitilde{}}\NormalTok{ physical\_3plus) }\SpecialCharTok{\%\textgreater{}\%}
  \FunctionTok{generate}\NormalTok{(}\AttributeTok{reps =} \DecValTok{1000}\NormalTok{, }\AttributeTok{type =} \StringTok{"bootstrap"}\NormalTok{) }\SpecialCharTok{\%\textgreater{}\%}
  \FunctionTok{calculate}\NormalTok{(}\AttributeTok{stat =} \StringTok{"diff in means"}\NormalTok{, }\AttributeTok{order =} \FunctionTok{c}\NormalTok{(}\StringTok{"yes"}\NormalTok{, }\StringTok{"no"}\NormalTok{)) }\SpecialCharTok{\%\textgreater{}\%}
  \FunctionTok{get\_confidence\_interval}\NormalTok{(}\AttributeTok{level =} \FloatTok{0.95}\NormalTok{)}

\NormalTok{weight\_ci}
\end{Highlighting}
\end{Shaded}

\begin{verbatim}
## # A tibble: 1 x 2
##   lower_ci upper_ci
##      <dbl>    <dbl>
## 1     1.11     2.41
\end{verbatim}

The 95\% confidence interval for the difference in mean weights between
students who are physically active at least three days per week and
those who are not is (1.11, 2.41) kg. We are 95\% confident that
students who are physically active at least three days per week have an
average weight between 1.11 kg and 2.41 kg higher than students who are
not physically active at least three days per week.

\subsection{More Practice}\label{more-practice}

\begin{enumerate}
\def\labelenumi{\arabic{enumi}.}
\setcounter{enumi}{7}
\tightlist
\item
  Calculate a 95\% confidence interval for the average height in meters
  (\texttt{height}) and interpret it in context.
\end{enumerate}

\textbf{Insert your answer here}

\begin{Shaded}
\begin{Highlighting}[]
\NormalTok{height\_ci }\OtherTok{\textless{}{-}}\NormalTok{ yrbss }\SpecialCharTok{\%\textgreater{}\%}
  \FunctionTok{drop\_na}\NormalTok{(height) }\SpecialCharTok{\%\textgreater{}\%}
  \FunctionTok{specify}\NormalTok{(}\AttributeTok{response =}\NormalTok{ height) }\SpecialCharTok{\%\textgreater{}\%}
  \FunctionTok{generate}\NormalTok{(}\AttributeTok{reps =} \DecValTok{1000}\NormalTok{, }\AttributeTok{type =} \StringTok{"bootstrap"}\NormalTok{) }\SpecialCharTok{\%\textgreater{}\%}
  \FunctionTok{calculate}\NormalTok{(}\AttributeTok{stat =} \StringTok{"mean"}\NormalTok{) }\SpecialCharTok{\%\textgreater{}\%}
  \FunctionTok{get\_confidence\_interval}\NormalTok{(}\AttributeTok{level =} \FloatTok{0.95}\NormalTok{)}

\NormalTok{height\_ci}
\end{Highlighting}
\end{Shaded}

\begin{verbatim}
## # A tibble: 1 x 2
##   lower_ci upper_ci
##      <dbl>    <dbl>
## 1     1.69     1.69
\end{verbatim}

The 95\% confidence interval for the average height of high school
students is 1.69 to 1.69 meters. We are 95\% confident that the true
average height of high school students represented by this sample falls
within this interval

\begin{enumerate}
\def\labelenumi{\arabic{enumi}.}
\setcounter{enumi}{8}
\tightlist
\item
  Calculate a new confidence interval for the same parameter at the 90\%
  confidence level. Comment on the width of this interval versus the one
  obtained in the previous exercise.
\end{enumerate}

\textbf{Insert your answer here}

\begin{Shaded}
\begin{Highlighting}[]
\NormalTok{height\_ci\_90 }\OtherTok{\textless{}{-}}\NormalTok{ yrbss }\SpecialCharTok{\%\textgreater{}\%}
  \FunctionTok{drop\_na}\NormalTok{(height) }\SpecialCharTok{\%\textgreater{}\%}
  \FunctionTok{specify}\NormalTok{(}\AttributeTok{response =}\NormalTok{ height) }\SpecialCharTok{\%\textgreater{}\%}
  \FunctionTok{generate}\NormalTok{(}\AttributeTok{reps =} \DecValTok{1000}\NormalTok{, }\AttributeTok{type =} \StringTok{"bootstrap"}\NormalTok{) }\SpecialCharTok{\%\textgreater{}\%}
  \FunctionTok{calculate}\NormalTok{(}\AttributeTok{stat =} \StringTok{"mean"}\NormalTok{) }\SpecialCharTok{\%\textgreater{}\%}
  \FunctionTok{get\_confidence\_interval}\NormalTok{(}\AttributeTok{level =} \FloatTok{0.90}\NormalTok{)}

\NormalTok{height\_ci\_90}
\end{Highlighting}
\end{Shaded}

\begin{verbatim}
## # A tibble: 1 x 2
##   lower_ci upper_ci
##      <dbl>    <dbl>
## 1     1.69     1.69
\end{verbatim}

\begin{Shaded}
\begin{Highlighting}[]
\NormalTok{width\_95 }\OtherTok{\textless{}{-}}\NormalTok{ height\_ci}\SpecialCharTok{$}\NormalTok{upper\_ci }\SpecialCharTok{{-}}\NormalTok{ height\_ci}\SpecialCharTok{$}\NormalTok{lower\_ci}
\NormalTok{width\_90 }\OtherTok{\textless{}{-}}\NormalTok{ height\_ci\_90}\SpecialCharTok{$}\NormalTok{upper\_ci }\SpecialCharTok{{-}}\NormalTok{ height\_ci\_90}\SpecialCharTok{$}\NormalTok{lower\_ci}

\NormalTok{width\_95}
\end{Highlighting}
\end{Shaded}

\begin{verbatim}
## [1] 0.003698088
\end{verbatim}

\begin{Shaded}
\begin{Highlighting}[]
\NormalTok{width\_90}
\end{Highlighting}
\end{Shaded}

\begin{verbatim}
## [1] 0.003082161
\end{verbatim}

The 90\% confidence interval is narrower (0.00299 meters) than the 95\%
confidence interval (0.00361 meters) 10. Conduct a hypothesis test
evaluating whether the average height is different for those who
exercise at least three times a week and those who don't.

\textbf{Insert your answer here}

\begin{Shaded}
\begin{Highlighting}[]
\NormalTok{obs\_height\_diff }\OtherTok{\textless{}{-}}\NormalTok{ yrbss }\SpecialCharTok{\%\textgreater{}\%}
  \FunctionTok{drop\_na}\NormalTok{(height, physical\_3plus) }\SpecialCharTok{\%\textgreater{}\%}
  \FunctionTok{specify}\NormalTok{(height }\SpecialCharTok{\textasciitilde{}}\NormalTok{ physical\_3plus) }\SpecialCharTok{\%\textgreater{}\%}
  \FunctionTok{calculate}\NormalTok{(}\AttributeTok{stat =} \StringTok{"diff in means"}\NormalTok{, }\AttributeTok{order =} \FunctionTok{c}\NormalTok{(}\StringTok{"yes"}\NormalTok{, }\StringTok{"no"}\NormalTok{))}

\NormalTok{obs\_height\_diff}
\end{Highlighting}
\end{Shaded}

\begin{verbatim}
## Response: height (numeric)
## Explanatory: physical_3plus (factor)
## # A tibble: 1 x 1
##     stat
##    <dbl>
## 1 0.0376
\end{verbatim}

\begin{Shaded}
\begin{Highlighting}[]
\NormalTok{null\_height }\OtherTok{\textless{}{-}}\NormalTok{ yrbss }\SpecialCharTok{\%\textgreater{}\%}
  \FunctionTok{drop\_na}\NormalTok{(height, physical\_3plus) }\SpecialCharTok{\%\textgreater{}\%}
  \FunctionTok{specify}\NormalTok{(height }\SpecialCharTok{\textasciitilde{}}\NormalTok{ physical\_3plus) }\SpecialCharTok{\%\textgreater{}\%}
  \FunctionTok{hypothesize}\NormalTok{(}\AttributeTok{null =} \StringTok{"independence"}\NormalTok{) }\SpecialCharTok{\%\textgreater{}\%}
  \FunctionTok{generate}\NormalTok{(}\AttributeTok{reps =} \DecValTok{1000}\NormalTok{, }\AttributeTok{type =} \StringTok{"permute"}\NormalTok{) }\SpecialCharTok{\%\textgreater{}\%}
  \FunctionTok{calculate}\NormalTok{(}\AttributeTok{stat =} \StringTok{"diff in means"}\NormalTok{, }\AttributeTok{order =} \FunctionTok{c}\NormalTok{(}\StringTok{"yes"}\NormalTok{, }\StringTok{"no"}\NormalTok{))}
\end{Highlighting}
\end{Shaded}

\begin{Shaded}
\begin{Highlighting}[]
\NormalTok{null\_height }\SpecialCharTok{\%\textgreater{}\%}
  \FunctionTok{get\_p\_value}\NormalTok{(}
    \AttributeTok{obs\_stat =}\NormalTok{ obs\_height\_diff,}
    \AttributeTok{direction =} \StringTok{"two\_sided"}
\NormalTok{  )}
\end{Highlighting}
\end{Shaded}

\begin{verbatim}
## # A tibble: 1 x 1
##   p_value
##     <dbl>
## 1       0
\end{verbatim}

The observed difference in average height between students who exercise
at least three days per week and those who exercise fewer than three
days per week is 0.0376. The simulated p-value is 0, meaning none of the
1,000 simulated permutations produced a difference as extreme as the
observed difference. This provides strong evidence against the null
hypothesis. Therefore, we reject the null hypothesis and conclude that
there is a statistically significant difference in average height
between the two physical activity groups

\begin{enumerate}
\def\labelenumi{\arabic{enumi}.}
\setcounter{enumi}{10}
\tightlist
\item
  Now, a non-inference task: Determine the number of different options
  there are in the dataset for the \texttt{hours\_tv\_per\_school\_day}
  there are.
\end{enumerate}

\textbf{Insert your answer here}

\begin{Shaded}
\begin{Highlighting}[]
\NormalTok{yrbss }\SpecialCharTok{\%\textgreater{}\%}
  \FunctionTok{summarise}\NormalTok{(}\AttributeTok{number\_of\_options =} \FunctionTok{n\_distinct}\NormalTok{(hours\_tv\_per\_school\_day, }\AttributeTok{na.rm =} \ConstantTok{TRUE}\NormalTok{))}
\end{Highlighting}
\end{Shaded}

\begin{verbatim}
## # A tibble: 1 x 1
##   number_of_options
##               <int>
## 1                 7
\end{verbatim}

There are 7 different options in the dataset for the
\texttt{hours\_tv\_per\_school\_day}.

\begin{enumerate}
\def\labelenumi{\arabic{enumi}.}
\setcounter{enumi}{11}
\tightlist
\item
  Come up with a research question evaluating the relationship between
  height or weight and sleep. Formulate the question in a way that it
  can be answered using a hypothesis test and/or a confidence interval.
  Report the statistical results, and also provide an explanation in
  plain language. Be sure to check all assumptions, state your
  \(\alpha\) level, and conclude in context.
\end{enumerate}

\textbf{Insert your answer here}

\begin{center}\rule{0.5\linewidth}{0.5pt}\end{center}

\end{document}
